\chapter{Design Evaluation and Testing}

\section{Prototype Construction}

As stated in the project scope, under section 1.3, the construction of the joint design was restricted to critical subsystems required to validate its core operating functions. As a result, instead of building the full-scale joint assembly, a scaled-down prototype of the primary electromagnetic brake was constructed.

Due to project constraints, inculuding an extended theorectical design phase, material availability and manufacturing facilities of the mechanical workshop, the physical build was adapted to the align with these capabilities.

\subsection{Mechanical Design of Prototype}

The resulting design implementation of the prototype resulted in a four-solenoid configuration, wound with copper wire with a \SI{0.28}{mm} diameter. This specific wire has a significanly reduced maximum current limit compared to the copper wire required for the \SI{6}{A} rating specified in the theoretcial design.

\subsection{Electrical Design of Prototype}

% Discuss the physical wiring of the switching circuit. 
% Detail the 24 V supply routing, the XL7015 buck converter stability, and the physical placement of the IRF640N MOSFET and freewheeling diode.

\subsection{Software Design of Prototype}

% Discuss the process of flashing the STM32 Nucleo. 
% Evaluate the real-world responsiveness of the TIM2 hardware timer tracking the quadrature pulses and the EXTI hardware interrupt.

\section{Safety and Performance Testing}

% Introduce the experimental testing protocols detailed in the safety report.

\subsection{Holding Force Verification}

% Detail the load testing methodology used to verify that the electromagnet successfully achieves the 128 N holding force requirement across the joint without slipping.

\subsection{Electrical Safety and Thermal Limits}

% Detail the physical multimetre measurements verifying the 6 A operating current. 
% Discuss the thermal dissipation of the MOSFET and the effectiveness of the freewheeling diode in mitigating back-EMF during physical switching events.